\section{Restricted class history and productivity sensitivity}
\label{a2:models}

The class-history candidate is a restricted two-class, shared-kernel
construction motivated by multivariate representations of marked excitation
\citep{DavisBaeumerWang2026,DavisKresinBaeumerWang2026}.
The equations below specify the implemented follow-up family; they are not
the full fitted model or a reproduction of the empirical comparisons in those
papers. It simplifies class and temporal-kernel structure while changing the
mark support, and is not a strict submodel of the authors' bounded-mark
family. Shared target kernels and the low-dimensional association parameter
are choices of this experiment. Notation is local: our \(\beta\) denotes the
magnitude rate, rather than a source paper's fractional-kernel order.

\subsection{Class association and the conditional ground intensity}

Write \(q=M-5\), let the class cut be \(c=0.5\), and define
\(\mathcal C_0=[0,c)\), \(\mathcal C_1=[c,\infty)\). Write
\(p_0=1-e^{-\beta c}\), \(p_1=e^{-\beta c}\) for the background class
probabilities under the working excess-exponential law. The common
\(\beta>0\) is estimated using training marks only and held fixed in each
forecast window. Let \(C_k\in\{0,1\}\) be the class label of event \(k\).
For source class \(j\), define the strictly past history
\[
 H_j(t)=\sum_{k:\,t_k<t,\ C_k=j}\kappa_{0.7}(q_k)g(t-t_k),
 \qquad
 R_0(\delta)=e^{-\delta/2},\quad R_1(\delta)=e^{\delta/2},
\]
where \(g\) is the retained unit-integral finite positive exponential
mixture and \(\kappa_r\) is defined in equation~\eqref{a2:productivity}.
Write \(\mu(t)=\exp\{a+d x(t)\}\) for the background, with the affine
observation coordinate specified in Section~\ref{a2:protocol}.
The implemented target-class intensity is
\begin{equation}
 \begin{aligned}
 \lambda_i(t)&=p_i\mu(t)+\frac{n}{\rho(h,\delta)}
       \sum_{j=0}^1 P_{ij}(h)R_j(\delta)H_j(t),\\
 P_{ij}(h)&=(1-h)p_i+h\,\mathbf 1_{\{i=j\}}.
 \end{aligned}
 \label{a2:class-intensity}
\end{equation}
Here \(n\) is the class-history candidate's scalar branching parameter and
\(\rho(h,\delta)\) is the normalizing Perron root specified below.
The frozen parameter box is
\[
 -0.95\min(p_0/p_1,p_1/p_0)\leq h\leq0.95,\qquad
 -1\leq\delta\leq1,\qquad 0\leq n\leq0.98.
\]
The factor \(0.95\) retains a margin inside the mathematical
nonnegativity limit for \(P\); the Perron root \(\rho(h,\delta)\) is
strictly positive.
For the working within-class magnitude laws, let
\[
 m_j=E[\kappa_{0.7}(Q)\mid Q\in\mathcal C_j],\qquad
 m_0=\frac{1-e^{-0.3\beta c}}{p_0},\quad m_1=e^{0.7\beta c}.
\]
The raw mean offspring matrix and its normalization are
\[
 A_{ij}(h,\delta)=P_{ij}(h)R_j(\delta)m_j,\qquad
 \rho(h,\delta)=\rho_{\rm PF}(A),\qquad K=\frac{n}{\rho}A,
\]
where \(\rho_{\rm PF}\) denotes the Perron root. Thus \(K\) is the
class-history candidate's mean offspring matrix, with Perron radius \(n\).
Association \(h\) and fertility contrast \(\delta\) play different roles.
Signed association can describe class anti-persistence while all
excitation coefficients remain nonnegative. At \(n=0\),
\(\lambda_i=p_i\mu\); trigger association and fertility are unidentified.

For every source class, \(\sum_iP_{ij}(h)=1\). Thus, conditional on a
fixed observed marked history,
\begin{equation}
 \Lambda(t)=\sum_i\lambda_i(t)
 =\mu(t)+s\{R_0(\delta)H_0(t)+R_1(\delta)H_1(t)\},
 \qquad s=\frac{n}{\rho(h,\delta)}.
 \label{a2:ground-intensity}
\end{equation}
The direct redistribution across target classes cancels from the ground
intensity. The conditional timing log likelihood is
\[
 \ell_{\rm time}=\sum_k\log\Lambda(t_k)-\int_{\mathcal W}\Lambda(t)\,dt.
\]
For these classes, the conditional within-class magnitude density is
\[
 f_i(q)=\frac{\beta e^{-\beta q}}{p_i}\mathbf1_{\{q\in\mathcal C_i\}}.
\]
The mark density at an observed event of class \(i\) is
\(\{\lambda_i(t)/\Lambda(t)\} f_i(q)\), giving mark log density
\(\log(\lambda_i/\Lambda)+\log\beta-\beta q-\log p_i\).
This gives the saved decomposition
\(\ell_{\rm joint}=\ell_{\rm time}+\ell_{\rm mark}\).

Equation~\eqref{a2:ground-intensity} has several qualifications.
First, holding \(n\) fixed does not remove \(h\), because \(\rho\) may
depend on \(h\). Second, matching an effective scale at \(h=0\) requires
\[
 n_0=n\frac{\rho(0,\delta)}{\rho(h,\delta)}.
\]
This mapping preserves timing intensities only if the mapped parameter is
admissible, in particular \(0\leq n_0\leq0.98\); \(n=0\) gives the
trivial no-trigger case. The frozen bound on \(n\), and any relevant
stability constraint, can make feasible effective-scale ranges depend on
\(h\). We therefore do not assert that all constrained timing families are
identical. Third, joint timing-and-mark estimation can shift \(s,\delta\)
and background parameters, producing indirect changes in timing forecasts.
Finally, the identity conditions on the actually observed marked past.
When future marks themselves are generated by the model, their class
composition can depend on \(h\), and so can a multistep or unconditional
timing law.

At \(h=0\), \(\lambda_i(t)=p_i\Lambda(t)\); the conditional magnitude law
is history-free even when \(\delta\) remains fitted. Consequently, we denote this association-zero ablation \(E1_{h=0}\);
it is the appropriate primary ablation for the association parameter
within this restricted family. It is not the same restriction
as setting both \(h\) and \(\delta\) to zero.

\subsection{Continuous productivity and dynamic evaluation}

The second follow-up replaces fixed \(r=0.7\) by a training-estimated
\(r\in[0,0.95]\) in
\begin{equation}
 \kappa_r(q)=(1-r)e^{\beta rq},\qquad
 \partial_r\kappa_r(q)=\kappa_r(q)
            \left(\beta q-\frac{1}{1-r}\right).
 \label{a2:productivity}
\end{equation}
For \(Q\sim{\rm Exp}(\beta)\), \(E[\kappa_r(Q)]=1\) when \(r<1\).
This is a normalization under the working law, not an empirical statement
about the historical magnitude catalogue.

The geometry records \(D_{ik}=g(t_i-t_k)\mathbf 1_{\{t_k<t_i\}}\)
and each parent's full-window exposure \(J_k\), independently of \(r\).
Every C++ value and gradient evaluation recomputes
\[
 H_r(t_i)=\sum_k D_{ik}\kappa_r(q_k),\quad
 K(r)=\sum_kJ_k\kappa_r(q_k)
\]
and their \(r\)-derivatives. In the excitatory candidate, and in the
inhibitory candidate whose supplied stress feature is independent of \(r\),
\begin{equation}
 \partial_r{\rm NLL}
 =n\left\{K'(r)-\sum_i\frac{\partial_rH_r(t_i)}{\Lambda(t_i)}\right\}.
 \label{a2:r-gradient}
\end{equation}
The common training magnitude law gives exactly the same mark score to
the fixed- and free-\(r\) stage-one candidates. Their score differences
therefore concern timing. At \(n=0\), the timing score is independent of
\(r\), its derivative vanishes and \(r\) is unidentified. Complete quiet
exposure, pretraining parents,
strictly past event histories and the final censor are retained.

For the same unbounded exponential working law and \(m>0\),
\[
 E[\kappa_r(Q)^m]=\frac{(1-r)^m}{1-mr}
 \quad\text{when }mr<1;
\]
the moment diverges for \(mr\geq1\). In particular, productivity variance
is infinite at \(r\geq0.5\), including the baseline \(r=0.7\).
Finite saved-catalogue likelihoods and numerical checks do not supply
variance-based confidence or power conclusions.

\subsection{Synthetic normalization and information}

The engineering controls are execution witnesses. For a multitype mean
offspring matrix \(K\), fixing its Perron radius does not fix the expected
first-generation offspring count of an immigrant with class law \(p\):
the latter is \(\mathbf 1^\top Kp\). The class-active history control
changed this quantity and the resulting offspring/event composition despite
using the same Perron branching parameter as its control. Its offspring
dilution must therefore be retained when interpreting successful fits or
forecast gains. Neither these controls nor the small fixed-\(r\) and
free-\(r\) engineering examples constitute matched-power, bias, coverage or
reliable-recovery studies.
